\section{From zero locations to the exterior interaction}\label{sec:exterior}

At a time \(T\) with a known all-real divisor, let \(x_-=x_k\) and
\(x_+=x_{k+1}\) be consecutive simple positive zeros. Write
\(g_k=x_+-x_-\) and \(D_k=g_k^2/4\). For integers \(p,q\ge0\) with
\(m=p+q\ge2\), the parent hierarchy uses
\[
 S_{k,p,q}=\sum_{\rho\ {\rm exterior}}
       (\rho-x_+)^{-p}(\rho-x_-)^{-q}.
\]
The sum counts multiplicities and omits just the two marked occurrences.
Thus the denominators are the displacements from the midpoint shifted by
the two signs of \(\sqrt{D_k}\). In the parent's notation
\(S_{k,p,q}=P^{(k)}_{p,q}\). On a separated real path the time derivative
of each entry uses the gap and finitely many entries of the same array,
including higher orders \cite[\S3, integer split recurrence]{AGC}.
For example, suppressing \(k\) on the right and left, an axis entry satisfies
\[
 S_{2,0}'=2\left(S_{2,0}^2-5S_{4,0}-2S_{3,1}+\frac{2}{g_k}S_{3,0}\right).
\]
The array is closed under differentiation when all orders and the gap are retained.
At a simple zero \(x_i\), differentiating \(H(t,x_i(t))=0\) gives
\(\dot x_i=H_{zz}/H_z\). Deflate its linear factor and differentiate
the paired product \eqref{eq:all-real-product}. This gives twice the
sum of \((x_i-\rho)^{-1}\) over the remaining zeros, paired by reflection.
In the difference of the two velocities the marked endpoints contribute
\(4/g_k\). Each common exterior occurrence contributes
\(-2g_k/((\rho-x_+)(\rho-x_-))\), and this difference is absolutely
summable. Consequently
\begin{equation}\label{eq:exterior-gap}
 M_k=S_{k,1,1},\qquad
 g_k'=\frac{4}{g_k}-2g_k M_k,\qquad D_k'=2-g_k^2M_k .
\end{equation}

Write \(Q_\rho=(\rho-c)^2-D=(\rho-x_-)(\rho-x_+)\) and
\(S_j=S_{k,j,j}=\sum_{\rho\ {\rm exterior}}Q_\rho^{-j}\) for the
diagonal traces. Thus \(M=S_1\) when the pair index is suppressed.
Differentiating the first mixed entry by the parent hierarchy
\cite[\S3, first diagonal row]{AGC} gives
\begin{equation}
 \dot D=2-4DM,\qquad
 \dot M=2M^2-\mathcal F,\qquad \mathcal F=12S_2+16D S_3 .
 \label{eq:parent}
\end{equation}
In particular the first derivative of the gap requires \(M\) and its
second derivative already involves higher exterior orders. We now give
a way to enclose these quantities from counted positions and source
moments at one time.

\paragraph{Completing the exterior from a counted head.}
Suppose the first \(K\ge k+1\) positive zeros are counted, and
\(x_+\le A<R\le x_i\) for every \(i>K\). Sum the contributions of the
first \(K\) reflected pairs directly, retaining \(-x_-\) and \(-x_+\).
For interval enclosures, each difference used with a positive exponent
must have a positive lower bound on its absolute value after allowing for
errors in both zeros. For the remaining pairs put
\[
 W_j=\sum_{i>K}x_i^{-2j},\qquad
 h_d=[v^d](1-x_+v)^{-p}(1-x_-v)^{-q}.
\]
The even product gives \(W_j\) by subtracting the counted head from the
full radial trace \(p_j\). In particular
\(W_1=\E_T U^2/2-\sum_{i\le K}x_i^{-2}\)
\cite{Founding07,EZMixedAnchor}. The complete tail of \(S_{k,p,q}\) equals
\begin{equation}\label{eq:exterior-tail}
 2\sum_{\substack{d\ge0\\m+d\ {\rm even}}} h_d W_{(m+d)/2}.
\end{equation}
Indeed reflection cancels terms of odd total degree. Since
\(0\le h_d\le \binom{m+d-1}{d} A^d\), for \(A/R<\sigma<1\)
the terms with \(d\ge L\) contribute at most
\[
  \frac{2R^{2-m} W_1}{(1-\sigma)^m}
       \left(\frac{A}{R\sigma}\right)^L .
\]
This bounds the truncation for any fixed row, including the axes.
An enclosure must also propagate the widths of the radial traces and
the head through their subtraction. For \(M_k\) the first tail term is
\(2W_1\).

\paragraph{A signed large-time readout.}
Set \(y_i=a x_i\), \(q_i=(i-\tfrac12)\pi\), and
\(\ell=a\sqrt{\kappa}\). Take \(K\asymp\ell\) inside a fixed multiple
of the large-time inverse window \cite{Founding08}. At fixed \(p,q\)
and \(0<\theta<1\), for \(k+1\le\theta K\) the counted separations and
the absolute tail bound give \cite{EZMixedAnchor}
\[
 a^{-m} S_{k,p,q}
 =\pi^{-m}\sum_{r\in\mathbb Z\setminus\{0,1\}}(r-1)^{-p}r^{-q}
    + O\!\left((\ell\sqrt{\kappa})^{-1}+K^{1-m}\right).
\]
For \(p=q=1\) the leading lattice interaction cancels the constant in
\eqref{eq:exterior-gap}. The more precise radial comparison of
\cite{EZMixedAnchor} gives, as \(T\to\infty\), uniformly for
\(k+1\le\sqrt{\ell}\),
\begin{equation}\label{eq:exterior-signed}
 \ell^2(g_k^2M_k-2)
   =\pi^2+O\!\left(\kappa^{-1/2}+\frac{(k+1)^2}{\ell}\right).
\end{equation}
The constants depend on the fixed window and the bounds on \(K/\ell\),
and in the general row estimate also on \(p,q,\theta\).
On the comparison lattice,
\(\sum_{r\in\mathbb Z\setminus\{0,1\}}1/[r(r-1)]=2\).
The correction that survives this cancellation can be read from the
first radial tail. Using \(\sum_{i\ge1}q_i^{-2}=1/2\) gives
\[
 2\sum_{i>K}(y_i^{-2}-q_i^{-2})
   =\frac{\Var_T U}{a^2}
      -2\sum_{i\le K}(y_i^{-2}-q_i^{-2}).
\]
To bound the three errors in this identity and the mixed sum, set
\(\beta=(\sqrt\kappa\,\ell^3)^{-1}\). On a disk of radius a fixed
multiple of \(\ell\), the inverse differs from the identity by
\(O(\kappa^{-1/2})\) and that difference vanishes to third order.
Cauchy's estimate on a smaller disk gives
\(|Y'(q)-1|\ll\beta |q|^2\). Extend the indices by
\(y_{1-j}=-y_j\). Integration between two levels now compares
their separation with its lattice value with relative error
\(O(\beta(q_i^2+q_j^2))\). In particular the positive radial head in
the last display contributes \(O(\beta K)\). The difference of the two
dimensionless mixed heads is bounded by
\[
 C\beta\!\!\sum_{\substack{1-K\le j\le K\\j\ne k,k+1}}
 \frac{j^2+(k+1)^2+1}{|j-k|\,|j-k-1|}
 \ll\beta(K+(k+1)^2).
\]
For the terms of positive degree in each reflected tail use
\(y_K\asymp K\), \(y_{k+1}=O(k+1)\), and
\(\sum_{j>K}y_j^{-2}=O(K^{-1})\). The last bound follows by
subtracting the counted head from the exact second trace as above,
using \(\sum_{j>K}q_j^{-2}=O(K^{-1})\) and \(\Var_T U/a^2=O(\ell^{-2})\).
The fourth-power tail is consequently \(O(K^{-3})\). Summing the even
degrees starting at two in \eqref{eq:exterior-tail} costs
\(O((k+1)^2 K^{-3})\) for each tail, since \(y_{k+1}/y_K<1/2\)
eventually. Finally \(a g_k=\pi(1+O(\beta(k+1)^2))\).
The only surviving term at scale \(\ell^{-2}\) for \(k+1=o(\sqrt\ell)\)
is \(\pi^2\Var_T U/a^2\). Substitution of
\(\kappa\Var_T U=1+O(\kappa^{-1/2})\) gives \eqref{eq:exterior-signed}.

Thus \(\ell^2(g_k^2M_k-2)\to\pi^2\) and \(D_k'<0\) uniformly on any
family \(k+1=o(\sqrt\ell)\) at sufficiently large \(T\). Section~\ref{sec:transport}
recovers the same leading drift on a wider large-time window by a phase estimate.
On a common large-time interval \([t,T]\), if \(t\to\infty\) and
\(\sup_{t\le s\le T}(k+1)^2/\ell_s\to0\), integration gives
\[
 D_k(t)-D_k(T)
     =(\pi^2+o(1))\int_t^T\frac{\dd s}{\ell_s^2}.
\]
This integration uses localization at each time in the interval. To obtain a time remainder from a single anchor we return to the field itself.

\subsection{The same data under backward evolution}
For \(s>0\) and \(P_m(z)=\prod_{i\le m}(1-b_i z^2)\), Gaussian convolution gives the exact earlier field
\begin{equation}
 \frac{H(T-s,z)}{Z(T)}
   =\E_T[e^{-sU^2}\cos(zU)]
   =\lim_{m\to\infty}\sum_{j=0}^m\frac{s^j}{j!}\partial_z^{2j}P_m(z).
 \label{eq:whole}
\end{equation}
Here \(b_i=x_i(T)^{-2}\) as in Section~\ref{sec:comparison}. Indeed \(|P_m(w)|\le H(T,\ii|w|)/Z(T)\); its log growth is \(O_T(1+|w|\log(2+|w|))\). Dominated convergence against the translated Gaussian therefore holds locally uniformly. The mass loss of the single frozen cosine with the exact initial curvature has the explicit size
\[
 \log\frac{Z(T-s)}{Z(T)e^{-s\E_T U^2}}
   =\int_0^s(s-v)\Var_{T-v}(U^2)\dd v .
\]
This follows by differentiating \(\log Z\) twice. Since \(\Var_T(U^2)\sim4L^2/\kappa\), the loss tends to one in logarithmic units at \(s=2\sqrt{T/L}\) \cite[\S4]{Founding07}. Equation~\eqref{eq:whole} gives a direct comparison for any backward estimate built from the same anchor data. Using a finite set of factors in that comparison requires a bound for the omitted product.
